\begin{center}
\textbf{\large{Tóm tắt}}
\end{center}

\addcontentsline{toc}{chapter}{Tóm tắt}

\begin{small}

Báo cáo trình bày quá trình phân tích, thiết kế và xây dựng ứng dụng di động \textbf{Fitty} nhằm hỗ trợ người dùng tập luyện thể hình và theo dõi sức khỏe cá nhân hóa trên nền tảng Android. Ý tưởng trung tâm của hệ thống là xây dựng một trợ lý số dành cho người tập gym, không chỉ cung cấp danh sách bài tập mà còn hỗ trợ người dùng chọn bài tập phù hợp, học đúng tư thế và động tác cơ bản, thiết lập kế hoạch tập luyện theo mục tiêu, duy trì lịch tập đều đặn, theo dõi sự thay đổi của cơ thể và hỗ trợ ghi nhận dinh dưỡng hằng ngày.

Để giải quyết bài toán, hệ thống được xây dựng theo kiến trúc nhiều tầng với Kotlin và Jetpack Compose ở tầng giao diện, lớp nghiệp vụ ở tầng điều phối, Room và DataStore ở tầng lưu trữ cục bộ, cùng Firebase làm hạ tầng phía sau cho xác thực, lưu trữ dữ liệu đám mây, lưu trữ media và thông báo đẩy. Một điểm kỹ thuật nổi bật của hệ thống là việc áp dụng mô hình \textit{offline-first} cho phân hệ bài tập, trong đó dữ liệu mô tả bài tập được đồng bộ từ Firestore, chuẩn hóa và lưu vào bộ nhớ cục bộ để giao diện có thể truy cập nhanh và ổn định hơn. Ngoài ra, hệ thống còn hỗ trợ quản lý nội dung động, cho phép cập nhật danh mục nhóm cơ, nội dung onboarding, mẫu kế hoạch tập khởi đầu và một số cấu hình hành vi mà không cần thay đổi trực tiếp phần mã giao diện.

Kết quả đạt được của báo cáo là một nguyên mẫu chức năng tương đối đầy đủ cho ứng dụng hỗ trợ gymer theo hướng cá nhân hóa. Hệ thống đã hiện thực được các phân hệ chính gồm: đăng ký và đăng nhập, dùng thử ở chế độ khách, khởi tạo hồ sơ tập luyện, xem trước và kích hoạt kế hoạch tập, quản lý lịch buổi tập, tra cứu bài tập theo nhóm cơ, xem chi tiết bài tập với mô tả, bước thực hiện, media hướng dẫn và gợi ý mức tập, thực hiện buổi tập nhanh hoặc buổi tập theo lịch, theo dõi các chỉ số như cân nặng, BMI, tỷ lệ mỡ cơ thể ước lượng, nhật ký bữa ăn và tiến độ theo ngày, đồng thời hỗ trợ nhắc việc và thông báo đẩy. Bên cạnh đó, dự án cũng đã có một số kiểm thử đơn vị cho các thành phần nghiệp vụ quan trọng như hình thành kế hoạch khởi đầu, đồng bộ dữ liệu bài tập, hoàn thành buổi tập và theo dõi tiến độ.

Nhìn chung, báo cáo cho thấy việc kết hợp kiến trúc Android hiện đại với các dịch vụ đám mây, lưu trữ cục bộ và các thành phần AI hỗ trợ là một hướng phù hợp để phát triển ứng dụng hỗ trợ tập gym có tính cá nhân hóa. Kết quả đạt được là nền tảng tốt để tiếp tục mở rộng hệ thống theo hướng nâng cao chất lượng gợi ý kế hoạch, cải thiện phân tích dinh dưỡng, trực quan hóa tiến độ cơ thể và hoàn thiện hơn trải nghiệm tập luyện thực tế.

\vspace*{1cm}
\textbf{Từ khóa}: Android, Jetpack Compose, Firebase, ngoại tuyến ưu tiên, ứng dụng tập luyện cá nhân hóa

\end{small}
